\documentclass[11pt]{article}

\usepackage[margin=1in]{geometry}
\usepackage[T1]{fontenc}
\usepackage{lmodern}
\usepackage{amsmath}
\usepackage{booktabs}
\usepackage{enumitem}
\usepackage{hyperref}
\usepackage{xcolor}
\usepackage{listings}

\hypersetup{
    colorlinks=true,
    linkcolor=blue,
    urlcolor=blue,
    citecolor=blue
}

\lstset{
    basicstyle=\ttfamily\small,
    breaklines=true,
    frame=single,
    columns=fullflexible,
    showstringspaces=false
}
\title{ADR-002\\Compile-Time Environment Policy Architecture}
\author{William Gordon Carter}
\date{\today}

\begin{document}

\maketitle

\begin{center}
\begin{tabular}{ll}
\toprule
\textbf{Field} & \textbf{Value} \\
\midrule
ADR & 002 \\
Title & Compile-Time Environment Policy Architecture \\
Status & Proposed \\
Author & William Gordon Carter \\
Created & \today \\
Supersedes & Original ADR-001 environment draft \\
Superseded By & --- \\
Depends On & ADR-001 Compile-Time Policy Architecture \\
\bottomrule
\end{tabular}
\end{center}

\tableofcontents
\newpage

\section{Context}

NavKit initially used Earth-specific terminology and constants, including
WGS-84, ECEF, and Earth gravity. This was appropriate for the initial prototype
but would tightly couple generic navigation algorithms to Earth.

The long-term NavKit architecture should support Earth, Moon, Mars, and other
central bodies. It should also support multiple gravity approximations,
mechanization frames, and future environment policies such as atmosphere and
magnetic field models.

This ADR applies the compile-time policy architecture defined in ADR-001 to
physical environment definitions.

\section{Decision}

NavKit shall represent physical environment assumptions using composable
compile-time policies.

The initial environment policy families are:

\begin{itemize}
    \item Planet policies
    \item Gravity policies
    \item Reference frame types
    \item Propagation / mechanization policies
\end{itemize}

Planet policies define physical constants and associated frames. Gravity
policies define gravitational acceleration equations. Mechanization policies
consume planet and gravity policies to propagate navigation state. The
Navigator remains an orchestration layer.

\section{Design Drivers}

\begin{itemize}
    \item Avoid Earth-specific assumptions in generic algorithms
    \item Support multiple celestial bodies
    \item Support multiple gravity fidelity levels
    \item Keep planetary constants separate from gravity equations
    \item Tie reference frames to central bodies explicitly
    \item Avoid inheritance hierarchy explosion
    \item Preserve zero-overhead compile-time configuration
    \item Keep the Navigator independent of direct environment details
\end{itemize}

\section{Planet Policies}

A planet policy defines immutable properties of a central body. The minimal
planet capability requires only the quantities needed by the lowest-fidelity
algorithms.

Example minimal concept:

\begin{lstlisting}[language=C++]
template<typename T>
concept PlanetPolicy = requires
{
    T::mu_m3_s2;
    typename T::InertialFrame;
    typename T::FixedFrame;
};
\end{lstlisting}

Additional concepts describe additional physical capabilities:

\begin{lstlisting}[language=C++]
template<typename T>
concept RotatingPlanetPolicy =
    PlanetPolicy<T> &&
    requires
    {
        T::omega_rad_s;
    };

template<typename T>
concept EllipsoidPlanetPolicy =
    PlanetPolicy<T> &&
    requires
    {
        T::a_m;
        T::b_m;
    };

template<typename T>
concept J2PlanetPolicy =
    PlanetPolicy<T> &&
    requires
    {
        T::J2;
        T::a_m;
    };
\end{lstlisting}

This capability-based structure avoids a rigid inheritance taxonomy.

\section{Concrete Planet Policies}

Concrete planet policies aggregate the facts they support.

Example:

\begin{lstlisting}[language=C++]
struct Wgs84 : PlanetPolicyBase<Wgs84>
{
    using InertialFrame = frames::EarthCenteredInertial;
    using FixedFrame = frames::EarthCenteredEarthFixed;

    static inline constexpr double mu_m3_s2 = 3.986004418e14;
    static inline constexpr double omega_rad_s = 7.292115e-5;

    static inline constexpr double a_m = 6378137.0;
    static inline constexpr double b_m = 6356752.314245179;

    static inline constexpr double J2 = 1.08262668e-3;
};
\end{lstlisting}

This single policy may satisfy:

\begin{itemize}
    \item \texttt{PlanetPolicy}
    \item \texttt{RotatingPlanetPolicy}
    \item \texttt{EllipsoidPlanetPolicy}
    \item \texttt{J2PlanetPolicy}
\end{itemize}

without explicitly declaring that it implements any concept.

\section{Gravity Policies}

Gravity policies define equations, not planetary constants.

A gravity policy should constrain its planet template parameter using the
weakest concept that expresses its requirements.

Examples:

\begin{lstlisting}[language=C++]
template<PlanetPolicy Planet>
struct Spherical;

template<J2PlanetPolicy Planet>
struct J2;
\end{lstlisting}

This allows valid low-fidelity configurations even when a planet does not
provide high-fidelity fields.

For example:

\begin{lstlisting}[language=C++]
using Gravity = gravity::Spherical<planet::SimpleMoon>;
\end{lstlisting}

may compile if \texttt{SimpleMoon} provides only \texttt{mu\_m3\_s2} and frame
types, while:

\begin{lstlisting}[language=C++]
using Gravity = gravity::J2<planet::SimpleMoon>;
\end{lstlisting}

will fail at compile time if \texttt{SimpleMoon} does not satisfy
\texttt{J2PlanetPolicy}.

\section{Reference Frames}

The term \emph{body-fixed} is avoided because \(b\) conventionally denotes the
vehicle body frame in navigation notation.

NavKit uses generic central-body terminology:

\begin{itemize}
    \item Planet-Centered Inertial (PCI)
    \item Planet-Centered Planet-Fixed (PCPF)
\end{itemize}

Planet-specific frame types include:

\begin{itemize}
    \item Earth-Centered Inertial
    \item Earth-Centered Earth-Fixed
    \item Moon-Centered Inertial
    \item Moon-Centered Moon-Fixed
    \item Mars-Centered Inertial
    \item Mars-Centered Mars-Fixed
\end{itemize}

Each planet policy exposes its frame types:

\begin{lstlisting}[language=C++]
using InertialFrame = ...;
using FixedFrame = ...;
\end{lstlisting}

This allows generic algorithms to use \texttt{Planet::InertialFrame} and
\texttt{Planet::FixedFrame} without being Earth-specific.

\section{Mechanization and Propagation Policies}

Mechanization is the layer where environment policies should enter the
navigation stack.

Example:

\begin{lstlisting}[language=C++]
using Planet = planet::Wgs84;
using Gravity = gravity::J2<Planet>;
using Propagation = nav::MechanizationPcpf<Planet, Gravity>;
using Nav = Navigator<Filter, Sensors, Propagation, UpdatePolicy>;
\end{lstlisting}

The concrete \texttt{Navigator} type indirectly encodes planet and gravity
through the propagation policy. This is acceptable. The important point is
that \texttt{Navigator} depends on a \texttt{PropagationPolicy} interface, not
directly on every lower-level environment policy.

\section{Dependency Direction}

The intended dependency direction is:

\begin{verbatim}
Navigator
    |
Propagation Policy / Mechanization
    |
+---+----------------+
|                    |
Gravity Policy   Frame Utilities
    \              /
     \            /
      Planet Policy
\end{verbatim}

Higher-level algorithms should depend on lower-level policies only through the
nearest appropriate abstraction.

\section{Compatibility Behavior}

Policy compatibility is enforced by concepts.

A simple gravity policy should compile for any planet satisfying the minimal
\texttt{PlanetPolicy}. A higher-fidelity gravity policy should compile only
for planets satisfying the required capability concept.

This creates a useful rule:

\begin{quote}
Algorithms do not depend on named planets; they depend on physical
capabilities.
\end{quote}

\section{Why Not a Planet Inheritance Hierarchy?}

A rigid hierarchy such as:

\begin{verbatim}
Planet
  -> RotatingPlanet
       -> EllipsoidPlanet
            -> J2Planet
\end{verbatim}

is rejected because physical capabilities are composable. A central body might
be rotating but not modeled as an ellipsoid, spherical but rotating, ellipsoid
without a supplied J2 coefficient, or supported only by a spherical gravity
model.

Concepts provide independent capability predicates without forcing a taxonomy.

\section{Why Not Template Navigator Directly on Planet?}

The following form is rejected:

\begin{lstlisting}[language=C++]
Navigator<Filter, Sensors, Planet, Gravity, Mechanization, UpdatePolicy>
\end{lstlisting}

This makes \texttt{Navigator} conceptually responsible for unrelated
environment details.

The preferred form is:

\begin{lstlisting}[language=C++]
Navigator<Filter, Sensors, Propagation, UpdatePolicy>
\end{lstlisting}

The propagation policy may internally encode planet, gravity, frame, and
mechanization choices. This preserves compile-time configuration while keeping
the high-level API organized.

\section{Current Implementation Direction}

The existing Earth-specific implementation should be refactored toward:

\begin{itemize}
    \item \texttt{planet/PlanetPolicy.hpp}
    \item \texttt{planet/PlanetPolicyBase.hpp}
    \item \texttt{planet/Wgs84.hpp}
    \item \texttt{gravity/GravityPolicy.hpp}
    \item \texttt{gravity/GravityPolicyBase.hpp}
    \item \texttt{gravity/Spherical.hpp}
    \item \texttt{gravity/J2.hpp}
    \item \texttt{frames/Frames.hpp}
\end{itemize}

Earth-specific compatibility aliases may be kept temporarily during the
transition, but new code should prefer policy-based naming.

\section{Alternatives Considered}

\subsection{Earth-Centric Implementation}

Rejected because it couples generic navigation code to Earth and ECEF.

\subsection{Single Planet Policy Owning Gravity}

Rejected because planetary constants and gravity equations are separate
concerns. The same planet may support multiple gravity policies.

\subsection{Runtime Environment Selection}

Rejected for the initial architecture because environment choices are static
for a given simulation, embedded build, or analysis campaign.

\subsection{Inheritance Taxonomy}

Rejected because independent capabilities compose more naturally than
hierarchical classes.

\section{Consequences}

\subsection{Benefits}

\begin{itemize}
    \item Earth assumptions become explicit and localized
    \item Moon and Mars support can be added without redesigning the estimator
    \item Gravity models can be swapped independently from planet definitions
    \item Frame ownership becomes type-level and planet-specific
    \item Mechanization becomes the primary abstraction for propagation
    \item Compile-time compatibility errors occur near the invalid policy use
\end{itemize}

\subsection{Costs}

\begin{itemize}
    \item More template types in the navigation layer
    \item More files and stronger naming discipline
    \item Requires understanding C++20 concepts
    \item Requires careful avoidance of unnecessary policy propagation upward
\end{itemize}

\section{Future Work}

Future environment policy families may include:

\begin{itemize}
    \item Atmosphere policies
    \item Magnetic field policies
    \item Earth orientation policies
    \item Geoid policies
    \item Terrain policies
    \item Higher-degree spherical harmonic gravity policies
\end{itemize}

These should follow ADR-001.

\section{Decision Summary}

NavKit shall treat the physical environment as a set of composable
compile-time policies. Planet policies aggregate physical constants and frame
types. Capability concepts describe which physical assumptions a planet
supports. Gravity policies constrain on the weakest planet capability required
by their equations. Mechanization policies consume planet and gravity policies
to propagate navigation state. The Navigator remains an orchestration object
parameterized by propagation behavior rather than by direct environment
details.

\end{document}
